\section{Theoretical Setup}

\subsection{Problem Setting}

Fix an integer $N\geq 1$, a support radius $R>0$, a density cap
$\kappa\in(0,\infty)$, and a compact interval $\Theta\subset\mathbb R$.
Let
\[
b:\Theta\to\mathbb R,
\qquad
F=(F_1,\ldots,F_N):\Theta\to\mathbb R^N,
\qquad
\phi_\alpha(\theta)=b(\theta)+\langle\alpha,F(\theta)\rangle.
\]
We use $\|x\|_1=\sum_i|x_i|$, and $|I|$ denotes one-dimensional
Lebesgue measure.  Unless positive length is required explicitly, intervals
$I$ are used as their literal subsets and may be open, closed, half-open,
empty, or singletons; they are not silently replaced by their closures.
For a nondegenerate compact interval $\Theta=[\theta_-,\theta_+]$, $C^1$
regularity means ordinary $C^1$ regularity on the interior with the first
derivative extending continuously to $\Theta$; derivatives at $\theta_-$
and $\theta_+$ are the corresponding one-sided endpoint derivatives.  If
$\Theta$ is a singleton, the derivative value is part of the specified
$C^1$ presentation, equivalently of a compatible local $C^1$ extension.
Put
\[
A=(2R)^N\kappa.
\]
Let $\mathcal D_{N,R,\kappa}$ be the class of Borel probability laws $\mu$
on $\mathbb R^N$ that have a full joint Lebesgue density $f_\mu$ satisfying
\[
f_\mu=0\quad\text{almost everywhere outside }[-R,R]^N,
\qquad
\|f_\mu\|_{L^\infty(\mathbb R^N)}\leq\kappa.
\]
This definition imposes no product structure or coordinate independence.

The deterministic family has a common one-dimensional triangular Pfaffian
presentation
\[
\eta_j'(\theta)=P_j(\theta,\eta_1(\theta),\ldots,\eta_j(\theta)),
\quad 1\leq j\leq q,
\]
\[
b(\theta)=Q_0(\theta,\eta(\theta)),
\qquad
F_i(\theta)=Q_i(\theta,\eta(\theta)),\quad 1\leq i\leq N.
\]
In ambient parameter dimension $p=1$, its descriptors are
\[
M=\max_{1\leq j\leq q}\deg P_j,
\qquad
\Delta_{\rm rnd}=\max_{1\leq i\leq N}\deg Q_i,
\qquad
\Delta_{\rm aff}=\max_{0\leq i\leq N}\deg Q_i.
\]
When $q=0$ the chain is absent and $M=0$.  The offset $b=Q_0$ contributes
to $\Delta_{\rm aff}$ but is not a random coefficient.

\subsection{Assumptions}

\begin{assumption}[Shared one-dimensional Pfaffian regularity]
\label{assump:shared-pfaffian-chain}
The functions $b,F_1,\ldots,F_N$ are $C^1$ on $\Theta$ and have the common
triangular Pfaffian representation and the
$(q,M,\Delta_{\rm rnd},\Delta_{\rm aff},p=1)$ degree convention above.
\end{assumption}

\begin{assumption}[Primitive no-forced-root nondegeneracy]
\label{assump:no-forced-root}
For every $\theta\in\Theta$,
\[
(b(\theta),F(\theta))\neq(0,0).
\]
Thus no parameter value is a root for every coefficient vector.  Points at
which $F(\theta)=0$ and $b(\theta)\neq0$ remain admissible and root-free.
\end{assumption}

\begin{assumption}[Arbitrarily correlated bounded joint density]
\label{assump:joint-density-cap}
The coefficient law is any $\mu\in\mathcal D_{N,R,\kappa}$.  The bound is
on the full joint density and does not impose independence.
\end{assumption}

\subsection{Objective and Scope}

Our objective is an unconditional affine coordinate-pivot sweep theorem.
It gives fixed-family finiteness of a static pivot-conditioning functional,
an ordinary-probability root bound uniform over admissible laws and
positive-length intervals, and two quantitative checks.  The first check
retains the exact $1/\delta$ scale in a rescaled two-feature family.  The
second keeps the leading coefficient of a monic polynomial deterministic
and recovers the exact coefficient for arbitrary correlated laws on its
lower coefficients.  No result below asserts a polynomial general-instance
bound for the conditioning functional in
$q,M,\Delta_{\rm rnd},\Delta_{\rm aff}$ or in other general Pfaffian
presentation data.
